% GENERATED FILE. Edit canonical lessons, then run npm run generate:books.
% canonical-source-hash: fnv1a64:be305bcb280ccfad
% canonical-lessons: JA-C519-omoinaosu, JA-C519-omoitodomaru, JA-C519-mikagiru, JA-C519-uragiru, JA-C519-bunan

\chapter{Giving Up and Thinking Again}
\label{ch:ja-giving-up-and-thinking-again}
\hlchaptermodality{519}

\begin{chapteropening}
\textbf{By the end of this chapter:} \emph{I can talk about giving up, thinking again and letting others down.}

Complete the last lesson of chapter 519: I can talk about giving up, thinking again and letting others down.
\end{chapteropening}

\section[\texorpdfstring{\ja{おもいなおす} (omoinaosu)}{おもいなおす (omoinaosu)}]{\ja{おもいなおす} (omoinaosu) — to think again; to change one's mind}

\label{lesson:JA-C519-omoinaosu}



\begin{quote}
Before the new one: say the Japanese for to slip up, then the Japanese for a fault.
\end{quote}

\subsection*{You'll want to know: \ja{おもいなおす}}
\textbf{\ja{おもいなおす}} — \emph{omoinaosu} — \textquotedblleft{}to think again; to change one's mind\textquotedblright{}.

\begin{quote}
\textbf{\ja{もし} \ja{いま} \ja{やめたら、かれは} \ja{あとで} \ja{おもいなおすでしょう。}} — \emph{moshi ima yametara, kare wa ato de omoinaosu deshō} — \textquotedblleft{}If he quits now, he will probably think again later.\textquotedblright{}
\end{quote}

Polite \textbf{\ja{おもいなおします}}, \emph{omoinaoshimasu}; a godan verb. Past (polite): \textbf{\ja{ちちが} \ja{ひっこしを} \ja{おもいなおしました。}}, \emph{chichi ga hikkoshi o omoinaoshimashita}.

\begin{cousinweb}
\textbf{\ja{おもいなおす}} (\emph{omoinaosu}) is a native Japanese word (wago), inherited rather than borrowed from Chinese, built from \emph{omou}, \textquotedblleft{}to think\textquotedblright{}, and \emph{naosu}, \textquotedblleft{}to fix\textquotedblright{}; its kanji was attached later as a kun'yomi (JMdict, KANJIDIC2).
\end{cousinweb}

\begin{grammarlens}[title={What you've built}]
One more word for saying what happens if.
\end{grammarlens}

\subsection*{Your turn}
\begin{itemize}
\raggedright
  \item \emph{Say it:} \emph{omoinaosu}
  \item \emph{Say it:} \emph{omoinaosu}, once more
  \item \emph{Say it:} \emph{moshi ima yametara, kare wa ato de omoinaosu deshō} — If he quits now, he will probably think again later.
  \item \emph{Recall:} say the Japanese for to slip up, then the Japanese for a fault, then say \emph{omoinaosu} again
\end{itemize}

\begin{tcolorbox}[breakable,colback=teal!4,colframe=teal!35!black,title={Before you move on}]
What does \ja{おもいなおす} mean? (\textquotedblleft{}To think again; to change one's mind\textquotedblright{}.) Say it once more.
\end{tcolorbox}
\section[\texorpdfstring{\ja{おもいとどまる} (omoitodomaru)}{おもいとどまる (omoitodomaru)}]{\ja{おもいとどまる} (omoitodomaru) — to think better of (doing)}

\label{lesson:JA-C519-omoitodomaru}



\begin{quote}
Before the new one: say the Japanese for a fault, then the Japanese for to think again.
\end{quote}

\subsection*{You'll want to know: \ja{おもいとどまる}}
\textbf{\ja{おもいとどまる}} — \emph{omoitodomaru} — \textquotedblleft{}to think better of (doing)\textquotedblright{}.

\begin{quote}
\textbf{\ja{もし} \ja{ははの} \ja{いけんを} \ja{きいたら、かれも} \ja{おもいとどまるでしょう。}} — \emph{moshi haha no iken o kiitara, kare mo omoitodomaru deshō} — \textquotedblleft{}If he listens to his mother's opinion, he too will probably think better of it.\textquotedblright{}
\end{quote}

Polite \textbf{\ja{おもいとどまります}}, \emph{omoitodomarimasu}; a godan verb. Past (polite): \textbf{\ja{あにが} \ja{ひっこしを} \ja{おもいとどまりました。}}, \emph{ani ga hikkoshi o omoitodomarimashita}.

\begin{cousinweb}
\textbf{\ja{おもいとどまる}} (\emph{omoitodomaru}) is a native Japanese word (wago), inherited rather than borrowed from Chinese, built from \emph{omou}, \textquotedblleft{}to think\textquotedblright{}, and \emph{todomaru}, \textquotedblleft{}to stay\textquotedblright{}; its kanji was attached later as a kun'yomi (JMdict, KANJIDIC2).
\end{cousinweb}

\begin{grammarlens}[title={What you've built}]
One more word for saying what happens if.
\end{grammarlens}

\subsection*{Your turn}
\begin{itemize}
\raggedright
  \item \emph{Say it:} \emph{omoitodomaru}
  \item \emph{Say it:} \emph{omoitodomaru}, once more
  \item \emph{Say it:} \emph{moshi haha no iken o kiitara, kare mo omoitodomaru deshō} — If he listens to his mother's opinion, he too will probably think better of it.
  \item \emph{Recall:} say the Japanese for a fault, then the Japanese for to think again, then say \emph{omoitodomaru} again
\end{itemize}

\begin{tcolorbox}[breakable,colback=teal!4,colframe=teal!35!black,title={Before you move on}]
What does \ja{おもいとどまる} mean? (\textquotedblleft{}To think better of (doing)\textquotedblright{}.) Say it once more.
\end{tcolorbox}
\section[\texorpdfstring{\ja{みかぎる} (mikagiru)}{みかぎる (mikagiru)}]{\ja{みかぎる} (mikagiru) — to give up on (someone)}

\label{lesson:JA-C519-mikagiru}



\begin{quote}
Before the new one: say the Japanese for to think again, then the Japanese for to think better of (doing).
\end{quote}

\subsection*{You'll want to know: \ja{みかぎる}}
\textbf{\ja{みかぎる}} — \emph{mikagiru} — \textquotedblleft{}to give up on (someone)\textquotedblright{}.

\begin{quote}
\textbf{\ja{もし} \ja{かれが} \ja{いつも} \ja{おくれたら、みんな} \ja{かれを} \ja{みかぎるでしょう。}} — \emph{moshi kare ga itsumo okuretara, minna kare o mikagiru deshō} — \textquotedblleft{}If he is always late, everyone will give up on him.\textquotedblright{}
\end{quote}

Polite \textbf{\ja{みかぎります}}, \emph{mikagirimasu}; a godan verb. Past (polite): \textbf{\ja{きゃくが} \ja{みせを} \ja{みかぎりました。}}, \emph{kyaku ga mise o mikagirimashita}.

\begin{cousinweb}
\textbf{\ja{みかぎる}} (\emph{mikagiru}) is a native Japanese word (wago), inherited rather than borrowed from Chinese; its kanji was attached later as a kun'yomi, and its older origin is not settled (JMdict, KANJIDIC2).
\end{cousinweb}

\begin{grammarlens}[title={What you've built}]
One more word for saying what happens if.
\end{grammarlens}

\subsection*{Your turn}
\begin{itemize}
\raggedright
  \item \emph{Say it:} \emph{mikagiru}
  \item \emph{Say it:} \emph{mikagiru}, once more
  \item \emph{Say it:} \emph{moshi kare ga itsumo okuretara, minna kare o mikagiru deshō} — If he is always late, everyone will give up on him.
  \item \emph{Recall:} say the Japanese for to think again, then the Japanese for to think better of (doing), then say \emph{mikagiru} again
\end{itemize}

\begin{tcolorbox}[breakable,colback=teal!4,colframe=teal!35!black,title={Before you move on}]
What does \ja{みかぎる} mean? (\textquotedblleft{}To give up on (someone)\textquotedblright{}.) Say it once more.
\end{tcolorbox}
\section[\texorpdfstring{\ja{うらぎる} (uragiru)}{うらぎる (uragiru)}]{\ja{うらぎる} (uragiru) — to betray; to let down}

\label{lesson:JA-C519-uragiru}



\begin{quote}
Before the new one: say the Japanese for to think better of (doing), then the Japanese for to give up on (someone).
\end{quote}

\subsection*{You'll want to know: \ja{うらぎる}}
\textbf{\ja{うらぎる}} — \emph{uragiru} — \textquotedblleft{}to betray; to let down\textquotedblright{}.

\begin{quote}
\textbf{\ja{もし} \ja{かれが} \ja{ともだちを} \ja{うらぎったら、ともだちは} \ja{おこるでしょう。}} — \emph{moshi kare ga tomodachi o uragittara, tomodachi wa okoru deshō} — \textquotedblleft{}If he betrays his friends, they will be angry.\textquotedblright{}
\end{quote}

Polite \textbf{\ja{うらぎります}}, \emph{uragirimasu}; a godan verb. Past (polite): \textbf{\ja{かれが} \ja{ともだちを} \ja{うらぎりました。}}, \emph{kare ga tomodachi o uragirimashita}.

\begin{cousinweb}
\textbf{\ja{うらぎる}} (\emph{uragiru}) is a native Japanese word (wago), inherited rather than borrowed from Chinese; its kanji was attached later as a kun'yomi, and its older origin is not settled (JMdict, KANJIDIC2).
\end{cousinweb}

\begin{grammarlens}[title={What you've built}]
One more word for saying what happens if.
\end{grammarlens}

\subsection*{Your turn}
\begin{itemize}
\raggedright
  \item \emph{Say it:} \emph{uragiru}
  \item \emph{Say it:} \emph{uragiru}, once more
  \item \emph{Say it:} \emph{moshi kare ga tomodachi o uragittara, tomodachi wa okoru deshō} — If he betrays his friends, they will be angry.
  \item \emph{Recall:} say the Japanese for to think better of (doing), then the Japanese for to give up on (someone), then say \emph{uragiru} again
\end{itemize}

\begin{tcolorbox}[breakable,colback=teal!4,colframe=teal!35!black,title={Before you move on}]
What does \ja{うらぎる} mean? (\textquotedblleft{}To betray; to let down\textquotedblright{}.) Say it once more.
\end{tcolorbox}
\section[\texorpdfstring{\ja{ぶなん} (bunan)}{ぶなん (bunan)}]{\ja{ぶなん} (bunan) — safe (as a choice); acceptable}

\label{lesson:JA-C519-bunan}



\begin{quote}
Before the new one: say the Japanese for to give up on (someone), then the Japanese for to betray.
\end{quote}

\subsection*{You'll want to know: \ja{ぶなん}}
\textbf{\ja{ぶなん}} — \emph{bunan} — \textquotedblleft{}safe (as a choice); acceptable\textquotedblright{}.

\begin{quote}
\textbf{\ja{もし} \ja{まよったら、この} \ja{みちが} \ja{ぶなんです。}} — \emph{moshi mayottara, kono michi ga bunan desu} — \textquotedblleft{}If you are unsure, this road is the safe choice.\textquotedblright{}
\end{quote}

A na-adjective: \ja{な} before a noun.

\begin{cousinweb}
\textbf{\ja{ぶなん}} (\emph{bunan}) is Sino-Japanese (kango): two Chinese characters read with their on'yomi, readings borrowed from Chinese. They are \emph{bu}, \textquotedblleft{}nothingness\textquotedblright{} (Middle Chinese \emph{mju}) and \emph{nan}, \textquotedblleft{}difficult\textquotedblright{} (\emph{nan}) (KANJIDIC2; Middle Chinese in Baxter's transcription).
\end{cousinweb}

\begin{grammarlens}[title={What you've built}]
One more word for saying what happens if.
\end{grammarlens}

\subsection*{Your turn}
\begin{itemize}
\raggedright
  \item \emph{Say it:} \emph{bunan}
  \item \emph{Say it:} \emph{bunan}, once more
  \item \emph{Say it:} \emph{moshi mayottara, kono michi ga bunan desu} — If you are unsure, this road is the safe choice.
  \item \emph{Recall:} say the Japanese for to give up on (someone), then the Japanese for to betray, then say \emph{bunan} again
\end{itemize}

\begin{tcolorbox}[breakable,colback=teal!4,colframe=teal!35!black,title={Before you move on}]
What does \ja{ぶなん} mean? (\textquotedblleft{}Safe (as a choice); acceptable\textquotedblright{}.) Say it once more.
\end{tcolorbox}
